% Peer reviews are normally anonymous; do not include your name.
% On a re-review, note the round in the title, e.g. "... (Round 2)".

\documentclass[11pt]{article}

% Shared typography for review PDFs: Times-style serif (newtx), 1in margins,
% and block paragraphs. Loaded by review.tex (\input) and by the Markdown build
% via pandoc (--include-in-header), and so a review produces the same-looking
% PDF whether it is written in Markdown or LaTeX.
\usepackage[margin=1in]{geometry}
\usepackage{newtxtext}
\usepackage{newtxmath}
\usepackage{parskip}
        % shared typography, also used by the Markdown build
\usepackage{enumitem}

\begin{document}

\section*{Review of: \emph{[Manuscript Title]}}

\subsection*{Summary of the submission}
% One or two paragraphs, in your own words: what the paper claims, what it
% does, and why it matters. This shows the editor you read it carefully.

\subsection*{Assessment}
% Your overall judgment and the single most important reason behind your
% recommendation. Leave the accept/reject decision itself to the editor.

\noindent\textbf{Recommendation:} [reject / major revision / minor revision / accept]

\subsection*{Major comments}
% Substantive issues affecting validity, novelty, soundness, or clarity of the
% contribution. Numbered so authors can respond point by point.
\begin{enumerate}
  \item
\end{enumerate}

\subsection*{Minor comments}
% Typos, wording, figure/table fixes, references.
\begin{enumerate}
  \item
\end{enumerate}

\subsection*{Questions for the authors}
\begin{enumerate}
  \item
\end{enumerate}

% ROUND 2+ : uncomment when re-reviewing a revision. Numbering refers to your
% round-1 review (see ../round-1/review.tex).
% \subsection*{Response to the revision}
% \begin{enumerate}
%   \item \textbf{[Prev major 1]} (resolved | partially addressed | not addressed):
% \end{enumerate}

\end{document}
